\documentclass[12pt, twoside]{report}

%Packages
\usepackage{setspace} %for onehalf spacing
\usepackage{graphicx} %for includegraphics
\usepackage{float} %for H pointer in graphics
\usepackage{titlesec} %to make section titles smaller
\usepackage{etoolbox} %remove page break in Bibliography
\usepackage{listings}
%Required packages in csfyp class:
\usepackage[top=1.2in, left=2.7cm, bottom=1.2in, right=2.7cm]{geometry} 
\usepackage{calc}
\usepackage[bottom]{footmisc}
\usepackage{caption}
\captionsetup[table]{skip=1ex}

%Parameters
\onehalfspacing
\renewcommand\thesection{\arabic{section}} %change numbering style of sections
%\titleformat*{\section}{\large\bfseries}
\titleformat*{\subsection}{\small\bfseries}
\patchcmd{\thebibliography}{\chapter*}{\section*}{}{}

\begin{document}
	\pagenumbering{roman}
	%Title Page:
	\begin{titlepage}
		\centering
		{\LARGE\bfseries Applied Cryptography: 2018/9 Assignment\par}
		\vspace{.5cm}
		
		{\Large \textbf{CPS3232} \par}
		\vspace{.5cm}
		
		{\large \textbf{Hands-on Crypto}\par}
		\vspace{0.5cm}
		
		{\large Marc Ferriggi (286397M)\par}
		\vfill
		
		\includegraphics[width=0.65\textwidth]{UoMLogo}\par
		\vfill
		
		{\large\bfseries Faculty of ICT \par}
		{\large\bfseries University of Malta \par}
		{\large\bfseries January 2019 \par}
	\end{titlepage}
	
	\tableofcontents
	\vfill
	\cleardoublepage
	\pagenumbering{arabic}
	
	\section{Block cipher modes and message authenticity}
	\pagebreak
	
	\section{Transparent access control on the Blockchain}
	\subsection{Implement the required reference monitor}
	\paragraph{ }The reference monitor was implemented in Jupyter Lab. The RSA modulus was generated using PyCryptodome's \textit{RSA.generate} function, where the number of bits was set to 2048 as this was deemed a sufficient length for 2017 and the public exponent \textit{e} was set to be 65537 as this is the FIPS standard. The generated key (which contains the RSA modulus $n=p.q$) was then saved to a file protected by a password in order to simulate that only Rene the reference monitor has access to it.
	
	\begin{lstlisting}[language=Python]
	%reset
	from Crypto.PublicKey import RSA
	import random
	\end{lstlisting}
	\begin{lstlisting}[language=Python]
	#Setting up FS
	#Generate n:
	key = RSA.generate(2048,e=65537) #nbits = 2048 (Sufficient length for 2017), e -> FIPS Standard
	#save the secure key to a file to only be accessed by Rene
	passphrase = 'Str0ngPassw0rd!%'
	f = open('mykey.pem','wb')
	f.write(key.exportKey('PEM',passphrase=passphrase)) #Use a passphrase since 
	f.close()
	n = key.n #assume n is public knowledge
	
	#Rene will then generate s, which lies in integer ring (Z_n)
	s = random.randint(1,key.n+1)
	#Assume this is sent to Alice and Bob over a secure channel
	
	#Delete Secrets:
	del key
	del passphrase 
	#Authorised parties will then be given mykey.pem and with knowledge of the passphrase will
	#have access to the secure key
	\end{lstlisting}
	
	\paragraph{ } Once the RSA key was set up, Rene then set \textit{v} by using the equation $v=s^2$ mod $n$.
	
	\begin{lstlisting}[language=Python]
	#Rene:
	#get the key
	f = open('mykey.pem','r')
	Rene_key = RSA.importKey(f.read(),passphrase='Str0ngPassw0rd!%')
	f.close()
	#set v
	v = (s**2)%Rene_key.n
	\end{lstlisting}
	
	\paragraph{ }The interactive protocol was then implemented setting $t=100$, thus having a successful attack probability of $2^{-100}$.
	
	\begin{lstlisting}[language=Python]
	t = 100
	for i in range(t):
		#Interactive Protocol
		#Alice
		#pick random r
		r = random.randint(1,n+1)
		x = (r**2)%n #compute x
		#x is now sent to Rene
		
		#Rene
		e = random.randint(0,1)
		#e is sent to Alice
		
		#Alice
		#if (i == 97):
		#    s = random.randint(1,n+1) #bad secret key
		y = (r*(s**e)) % n
		#y is sent to Rene
		
		#Rene:
		if (y**2 % n != (x*(v**e))%n):
			print('Attack!')
			break
	\end{lstlisting}
	
	\subsection{Attack 1}
	
	\paragraph{ }Weaken the implementation in a way so that Alice always commits to the same $r \in Z_n$. Demonstrate an attack that discloses $s$ and which hinges on the fact that whenever $e = 0$ Alice sends $y ≡ r$ $mod$ $n$ as a response.
	
	\paragraph{ }The implementation of the reference monitor was then weakened as specified. Given the fact that $r$ wasn't changing, Oscar could simply extract $r$ by intercepting $y$ when $e=0$. Once Oscar has $r$ then, when $e=1$, Oscar intercepts $y$ and can extract $s$ by using the equation $s=y.r^{-1}$ $mod$ $n$. The modified code to the previous implementation can be seen below.
	
	\begin{lstlisting}[language=Python]
	from sympy import mod_inverse
	
	gotR = 0
	
	t = 100
	#Alice will pick a random r and keep it fixed:
	r = random.randint(1,n+1)
	for i in range(t):
		#Interactive Protocol
		#Alice:
		x = (r**2)%n #compute x
		#x is now sent to Rene
		
		#Rene
		e = random.randint(0,1)
		#e is sent to Alice on open channel, Oscar can intercept this
		
		#Alice
		y = (r*(s**e)) % n
		#y is sent to Rene on an open channel, Oscar can intercept this
		
		#Oscar's Attack:
		if e==0 and gotR == 0:
			oscar_r = y
			gotR = 1
		if e==1 and gotR == 1:
			oscar_s=(y*mod_inverse(r,n)) % n
			print(s==oscar_s)
			break
	
		#Rene:
		if (y**2 % n != (x*(v**e))%n):
			print('Attack!')
			break
	\end{lstlisting}
	
	
	\pagebreak
	\section{Setting up Certificate Authority (CA) trees and HTTPS servers}
\end{document}